\section{Government, mission, and guardianship}

\subsection{Question 110: Spiritual agents and bodily things}\label{q110}
Aquinas places angelic activity within the government of creation. Article~1 argues from the relation of particular bodily powers to the more universal reach of an intellectual agent. His account includes the cosmology of his age, in which spiritual substances move the heavens. Yet he does not confine angelic ministry to celestial motions: the government of lower bodily things also belongs within providence. The nine orders describe general offices, not a separate choir for every bodily species.\footnote{\ST{I q.~110 a.~1}, especially ad~2--3.}

Article~2 asks a sharper question: does matter receive whatever form an angel simply wills? Aquinas answers no. A creature's superiority does not confer the Creator's sovereignty over being. Angels can bring bodily agents into action, but their willing is not itself an immediate source of substantial forms. Aquinas compares a cook who regulates a fire: the resulting dish depends on intelligent direction, yet it is produced through bodily processes. Augustine similarly distinguishes God's inward creative causality from the outward arrangement of created causes by angels or human beings.\footnote{\ST{I q.~110 a.~2}; Augustine, \emph{De Trinitate} III.8.13--9.17.}

The immediate bodily power affirmed in article~3 is local motion. An angel can move bodies and thereby bring natural agents into contact or arrange conditions for effects. This explains how the preceding denial can coexist with real angelic action: denying unrestricted command over matter is not denying every capacity to act upon it. The natural capacities employed remain created capacities.\footnote{\ST{I q.~110 a.~3}.}

Article~4 reserves a miracle strictly understood to God. A surprising event exceeds neither nature nor every created power simply because it exceeds our experience. For Aquinas, a true miracle exceeds the order of the whole created universe, which includes angelic powers. Angels may intercede for such a work or minister in its accomplishment; their natural power does not produce it as a principal cause. Accordingly, Gregory's attribution of wonders to the Virtues can be affirmed without making an angel a second divine wonder-worker.\footnote{\ST{I q.~110 a.~4}, especially ad~1--2 and ad~4; Gregory, \emph{Homilies on the Gospels} II.34.10.}

\angeldossier{I q.~110, aa.~1--4}{Thomist position: angelic government through created bodily causes, under the Creator's universal power.}{Augustine, \emph{De Trinitate} III.4,8--9; Gregory, \emph{Homilies} II.34.10; Psalm 135:4 (Hebrew 136:4).}{The account rejected in a.~2 makes a separated intellect the immediate source of material forms.}

\subsection{Question 111: Helping a mind without taking its place}\label{q111}
In article~1 angelic illumination respects human embodiment. An angel strengthens understanding and presents truth through sensible likenesses, because human thinking in this life turns to images. The angel does not implant the infused virtue of faith. Aquinas distinguishes proposing what is to be believed from the divine gift through which a person believes. The same distinction allows a human preacher to serve faith without creating faith by persuasion. Nor does receiving assistance entail knowing which angel, or even knowing that an angel, assisted.\footnote{\ST{I q.~111 a.~1}, corpus and ad~1--3.}

Article~2 repeats and extends the protection of freedom. Only God moves the will inwardly as its creator. An angel can present a good for consideration; with human beings, an angel can also affect passions through their bodily conditions. Neither mode makes the will's consent necessary. Fear, anger, or desire can incline a choice without constituting the choice. Aquinas's distinction applies equally to a good angel's help and a demon's solicitation: a created spirit cannot seize the will and turn its own intention into the person's voluntary consent.\footnote{\ST{I q.~111 a.~2}, corpus and ad~2--3.}

Articles~3 and~4 distinguish imagination from external sensation. Aquinas holds that good or bad angels can act upon bodily conditions of imagination and sense. He explains those conditions through the movement of bodily spirits and humors, the physiology he received; this is not a modern account of the nervous system. His conceptual limit is that an angel rearranges material furnished through sense rather than supplying a wholly new sensory capacity. His example is that an angel cannot make someone blind from birth imagine color. Sense can also be affected through external objects or through changes within the organ. The articles state powers in general; they do not identify the cause of a particular person's experience.\footnote{\ST{I q.~111 aa.~3--4}, especially a.~3 ad~1--3.}

\angeldossiersettled{I q.~111, aa.~1--4}{Thomist position: intellectual illumination and influence upon imagination, sense, and passion; God alone inwardly causes the act of willing and infuses faith.}{Dionysius, \emph{Celestial Hierarchy} I, IV; Proverbs 21:1; Matthew 1:20; Genesis 19:11.}

\subsection{Question 112: Sent without ceasing to contemplate}\label{q112}
To be sent is to receive a commission and apply one's power in its execution. In article~1 an angel's mission is a new exercise of limited created power under divine authority. Service to humanity does not lower the angel's dignity: the angel principally serves God, who wills that service. Neither does external ministry interrupt the vision of God. Unlike our divided attention, angelic contemplation supplies the rule for the action; the minister acts from the divine presence rather than taking leave of it.\footnote{\ST{I q.~112 a.~1}, especially ad~2--4.}

Article~2 asks whether every order is sent on external missions. Aquinas answers that the higher orders communicate through the lower, and denies exceptional direct external missions of those higher orders. Isaias's Seraph can accordingly be understood as a lower minister performing a seraphic work, or as the action of a higher angel effected through a lower. This is Aquinas's definite resolution. Gregory reports the Dionysian account but then declines to decide, without clear testimony, whether such higher spirits act directly or through subordinate ranks. Gregory's reserve is more extensive than Aquinas's conclusion.\footnote{\ST{I q.~112 a.~2}, corpus and ad~2; Dionysius, \emph{Celestial Hierarchy} XIII; Gregory, \emph{Homilies} II.34.12--13.}

Article~3 distinguishes two senses of standing before God. All blessed angels stand before him in the immediate vision of his essence, including those sent to us. In a narrower sense, the higher angels stand before him by receiving disclosures of divine mysteries which they communicate to others. An account of mediated knowledge must never become an account of a mediated beatific vision.\footnote{\ST{I q.~112 a.~3}.}

Article~4 identifies five orders with external execution: Virtues, Powers, Principalities, Archangels, and Angels. Dominations direct the ministry without themselves executing it externally; the three orders of the first hierarchy likewise remain above that execution. Thus six orders may be called ministering when direction is included, although only five execute external commissions. Aquinas also discusses different readings of Daniel's immense numbers. Those readings signify an abundance beyond human reckoning, not a disclosed numerical census.\footnote{\ST{I q.~112 a.~4}, corpus and ad~1--2.}

\angeldossier{I q.~112, aa.~1--4}{Church teaching: angelic ministry, Catechism nos.~331--333. Thomist position: the distribution of ministerial offices. Disputed: whether the higher orders ever execute external missions directly.}{Exodus 23:20; Hebrews 1:14; Daniel 7:10; Dionysius, \emph{Celestial Hierarchy} VIII, XIII; Gregory, \emph{Homilies} II.34.12--13.}{Gregory expressly leaves direct action by Cherubim and Seraphim unresolved. Aquinas rules out direct external mission by the higher orders and also rejects an unnamed opinion allowing exceptions.}

\subsection{Question 113: The guardian and the person}\label{q113}
Guardianship belongs to providence's care for free but fallible persons. In article~1 Aquinas argues that knowing a general moral principle does not guarantee applying it rightly in a particular circumstance. Angelic help addresses that weakness. God remains the giver of grace and virtue; angelic instruction assists the person's practical direction. A person can resist that assistance. Sin therefore does not demonstrate negligence in the guardian.\footnote{\ST{I q.~113 a.~1}, especially ad~1--3.}

Article~2 assigns an individual guardian to each person. Its reason is the enduring value of the individual rational soul, not merely the preservation of humanity as a species. A public guardian's care for a community does not replace care for each member's interior good. Article~3 assigns this personal office to the lowest order and distinguishes it from broader guardianship exercised through higher orders. Aquinas interprets Chrysostom's appeal to greater angels as compatible with greater members of the lowest order. Chrysostom himself emphasizes the dignity of the supposedly insignificant people under their care and does not specify that scholastic solution.\footnote{\ST{I q.~113 aa.~2--3}, especially a.~3 ad~1; John Chrysostom, \emph{Homilies on Matthew} 59.4, on Matthew 18:10.}

In article~4, the protection extends to every human wayfarer, including unbelievers and those who resist grace. Its benefits can include restraining harm even when the recipient does not attain salvation. Christ is the distinct case: angels minister to him, but he needs no guardian superior to himself. His human soul is united to the Word and possesses the vision of God even during his earthly journey. Aquinas also distinguishes a guardian's office during the journey from the communion of angels and redeemed humans at its end.\footnote{\ST{I q.~113 a.~4}, especially ad~1 and ad~3.}

Article~5 places assignment of an individual guardian at birth rather than baptism. The benefit belongs to providence for rational creatures; it is not limited to sacramental incorporation into Christ. Origen had posed birth and baptism as alternatives for examination, without supplying the simple categorical baptism-only doctrine sometimes attributed to him. Aquinas decides for birth, invoking Jerome. He judges it probable that the mother's angel guards the child in the womb.\footnote{\ST{I q.~113 a.~5}, corpus and ad~3; Origen, \emph{Commentary on Matthew} XIII.27--28; Catechism of the Catholic Church, no.~336.}

The guardian never abandons the person entirely, according to article~6. Providence can permit an affliction or a sinful act which the angel does not prevent. Such permission is not cancellation of the office. Article~7 adds that the blessed angels' care does not destroy their happiness. They do not approve sin; their will adheres to God's just ordering and permission, rather than becoming frustrated opposition to providence. The distinction concerns blessed participation in divine wisdom, not indifference to human suffering.\footnote{\ST{I q.~113 aa.~6--7}.}

Finally, article~8 interprets the resistance of the prince of Persia in Daniel 10:13. On the interpretation Aquinas develops, good angels representing contrary claims seek the divine judgment; their wills do not quarrel against one another or against God. The difficulty lies in the opposed matters presented for judgment and in what has not yet been disclosed to them. Even faithful ministers do not possess God's comprehensive knowledge of providence.\footnote{\ST{I q.~113 a.~8}.}

\angeldossier{I q.~113, aa.~1--8}{Church teaching: angelic care of human life, Catechism nos.~336, 352. Thomist position: the individual guardian's assignment, rank, and commencement, and the reading of Daniel's resistance.}{Psalm 90:11 (Hebrew 91:11); Matthew 18:10; Daniel 10:13; Apocalypse 21:4. Aquinas invokes Jerome, Origen, Chrysostom, and Gregory.}{Origen, \emph{Commentary on Matthew} XIII.27--28, explores birth and baptism as alternatives; Aquinas chooses birth. Chrysostom's greater angels do not specify a choir.}
